\section{Control Development}
\label{sec:ctrl-comparison}

Control development addressed manual voltage adjustment, blocking conversion waits, an unsuitable feedback quantity, and excess energy transfer during startup or after load removal (\Cref{tab:ctrl-version-identity}). Each stage changed how the switching bursts responded to the converter and its load.

\begin{table}[htbp]
\centering
\caption{Firmware changes and their functional improvements.}
\label{tab:ctrl-version-identity}
\begin{tabularx}{\textwidth}{@{}L{0.17\textwidth}L{0.37\textwidth}Y@{}}
\toprule
Stage & Change & Improvement\\
\midrule
Open loop & Manual carrier and burst-enable commands & Check switching timing and basic transfer before adding automatic voltage control\\
Blocking feedback & Voltage feedback with blocking ADC reads & Adjust transfer from the sensed voltage instead of relying on manual correction\\
DMA feedback & Continuous ADC DMA and corrected feedback scaling & Remove conversion waits and regulate bulk voltage from PA1 rather than an unsuitable voltage--current product\\
Final baseline & Burst-edge alignment attempts, startup precharge, no-load guard, and loaded duty trim & Limit commanded energy at startup and no load, reduce continued charging after load removal, and adjust transfer as load changes\\
V2 calibration and setup & Browser calibration, reference-load settings, and output-OFF carrier timing & Repeat duty searches for known loads and vary switching frequency and dead time without rebuilding the firmware\\
\bottomrule
\end{tabularx}
\end{table}

\subsection{Open-Loop Bring-Up}
\label{sec:ctrl-v0}

Manual carrier and enable commands were used to check timing, the dead interval, and basic energy transfer under controlled input and load. Each command could be changed directly, making this a useful mode for bringing up the power stage and finding connection or timing problems. Output voltage still varied with load, and continued bursts could charge the unloaded output, so closed-loop control was needed. Open-loop operation was kept for diagnostics.

\subsection{Voltage Control with Blocking ADC Reads}
\label{sec:ctrl-v1}

Voltage feedback allowed the burst command to increase or reduce energy transfer as the sensed voltage changed under load, replacing manual adjustment. The first implementation waited for each ADC conversion before updating the command. These conversion waits shared the foreground execution path with UART handling and other work, which could delay the next duty adjustment.

\label{sec:ctrl-blocking}
The conversion wait lengthened the time between opportunities to process a command or update burst duty. Circular DMA removed that wait by transferring samples into memory in the background. ADC acquisition could then continue while the foreground handled commands and control.

\subsection{Moving to DMA and Voltage Feedback}
\label{sec:ctrl-v2}

The baseline uses continuous three-channel ADC1 acquisition and circular DMA with the sample-time setting \texttt{ADC\_SAMPLETIME\_39CYCLES\_5}. Samples are transferred to memory in the background, so the foreground can read the latest available values without waiting in \texttt{analogRead()} while handling commands and control updates.

The feedback quantity also had to be corrected. In early tests, voltage was multiplied by the current-related input, but PA2 carries half-bridge-current information rather than connector output current. That product was removed from active regulation, and the loop was closed on secondary voltage. The error signal then corresponded to the bulk supply the MCU was intended to regulate. The \texttt{POWER} command, \texttt{KP}, \texttt{KI}, and earlier power-feedback fields remain for compatibility or diagnostics. They do not select an active PI output-power controller in the baseline firmware.

Startup precharge limits the initial energy packets used to power the isolated feedback circuit. The no-load guard and soft holdoff reduce continued charging after load removal, while loaded duty trim adjusts transfer as demand changes. Bounded carrier-update polling attempts to place burst edges at carrier boundaries to reduce partial cycles without changing the fixed carrier.

